% Section 99 --- References. Compiled from the endnotes and prose of sections
% 00--40; bibliographic detail, witness identifications, and rights notes per
% the research dossiers in build/ark-research/ (esp. w2-rights.md).
\sectionguard
\section{References}

Witnesses are listed with role and rights notes; the access date for all
web loci is 2026-08-15 unless another date is stated in the endnote where
the witness is used. Stable source identities are registered, or to be
registered, in the repository source library; this publication's bindings
are in \texttt{research/source-bindings.toml}. The list contains only
works cited in the study's endnotes or named in its prose. Entries give
the principal witnesses; subsidiary transcriptions and discussion pages
used for single points stand in the endnotes at their loci, with exact
URLs. Works consulted at one remove --- known through reviews, through
quotations transmitted verbatim in other works read for this study, or
through search-level reporting --- are listed as such; the body claims
nothing from them beyond what the endnotes attribute.

\subsection*{Scripture and versions}

\begin{itemize}
\item Douay-Rheims Bible (Challoner revision; public domain), drbo.org
chapter pages (e.g.\ \url{https://drbo.org/chapter/10006.htm} for
2 Kings 6). Role: every English Scripture quotation, with the Douay book
names (Josue, 1--4 Kings, 1--2 Paralipomenon, Jeremias, 2 Machabees,
Apocalypse) and Vulgate psalm numbers (77 = Hebrew 78; 131 = 132).
Challoner's note at 1 Kings 6:19 was read on the same chapter page.
\item Clementine Vulgate, as printed at drbo.org beside the English.
Role: the Latin of Exodus 25:17, 22 and 3 Kings 6:19. Limit: the Latin
of Ps 131:8 was not collated there; \term{arca sanctificationis tuae}
rests on the AAS and Migne witnesses below.
\item Septuagint research instruments, described, never base text: the
Elpenor edition (ellopos.net; versification offsets noted where cited);
the machine-readable Rahlfs 1935 corpus of E.~Wong,
\url{https://github.com/eliranwong/LXX-Rahlfs-1935}, on which the
lemma-level concordance searches were run and adversarially re-run; the
Blue Letter Bible LXX pages (blueletterbible.org) confirming the
decisive loci. A spot-check against the printed Rahlfs--Hanhart remains
open (section 90).
\item \work{SBL Greek New Testament},
\url{https://github.com/LogosBible/SBLGNT}. Role: NT concordance counts
and the critical text of Lk 1:42; Nestle 1904, Westcott--Hort, and the
Byzantine tradition cited comparatively there.
\item Revised Standard Version, read at biblegateway.com. Role:
comparative modern witness only --- its translators' note at 1 Samuel
6:19, its renderings where the Douay diverges, the description of
4 Esdras 10:21--22. In copyright; brief attributed use only.
\item Abegg, Flint, and Ulrich, \work{The Dead Sea Scrolls Bible}
(1999), read in an archive.org lending copy. Role: the 4QSam-a
evidence (coverage at 1 Kings 6; readings at 2 Kings 6:2, 6--7, 13). In
copyright; reported and paraphrased. DJD XVII (2005) not consulted.
\end{itemize}

\subsection*{Magisterial and liturgical}

\begin{itemize}
\item Pius XII, \work{Munificentissimus Deus} (1 November 1950), Vatican
English web text,
\url{https://www.vatican.va/content/pius-xii/en/apost_constitutions/documents/hf_p-xii_apc_19501101_munificentissimus-deus.html};
Latin at \work{AAS} 42 (1950) 753--773, read in the scan
\url{https://www.vatican.va/archive/aas/documents/AAS-42-1950-ocr.pdf}.
Role: the Assumption definition; the ark at \S\S 26, 29, 34 and there
only; the footnote-20 finding; also the witness for Anthony of Padua
(\S 29) and Bellarmine (\S 34), quoted only as transmitted there. Two
defects of the vatican.va English are recorded in the notes.
\item Pius IX, \work{Ineffabilis Deus} (1854), English at
papalencyclicals.net, translator unnamed; the Latin not checked against
the \work{Acta}. Role: ``the ark and house of holiness.''
\item Second Vatican Council: \work{Dei Verbum} (1965) \S 12,
paraphrased throughout, no edition quoted; \work{Lumen Gentium} (1964),
English and Latin from vatican.va, searched in full --- doctrinal frame
of the coda (LG 53, 60--62, paraphrased) and the verified negative for
\term{arca}.
\item Paul VI, \work{Marialis Cultus} (1974), English and Latin at
vatican.va. Role: \term{vera Foederis Arca} at MC 6; the ark clause and
curated patristic note at MC 26 (Severus of Antioch, Hesychius,
Chrysippus, Andrew of Crete, John Damascene; the Akathist in the
companion note); the negatives recorded with them.
\item \work{Catechism of the Catholic Church}: Latin editio typica at
\url{https://www.vatican.va/archive/catechism_lt/p4s1c2a2_lt.htm};
English second edition at the St.~Charles Borromeo mirror
(scborromeo.org). Role: \S\S 112--114, 128--130 (typology, paraphrased);
\S\S 489, 722, 2676--2677 (the Daughter-Zion and ark fusion).
\item John Paul II: Angelus of 10 July 1983, Italian at vatican.va
(document hf\_jp-ii\_ang\_19830710), quoted briefly with editorial
paraphrase; \work{Redemptoris Mater} (1987), searched --- verified
negative for ``ark''; general audience of 2 October 1996, Italian
(hf\_jp-ii\_aud\_19961002), searched --- no \term{arca}.
\item \work{Breviarium Romanum}, Pars Aestiva, 1942 printing,
\url{https://archive.org/details/breviarium-romanum-1942}. Public
domain. Role: the pre-1950 Matins of 15 August (the Damascene lesson
\term{Hodie sacra et animata arca}\ldots, quoted as printed; the
Fortunatus hymn \work{Quem terra, pontus, sidera} in the Urban VIII
revision) and the bound-in 1950 supplement (\term{Arca non putri
fabricata ligno}). Latin quoted with labeled editorial paraphrases.
\item Roman Lectionary (1969) and current order of readings, verified at
the USCCB readings pages, \url{https://bible.usccb.org} (Lectionary
nn.~621--622; the Presentation responsorial at the daily-readings page
for 2 February 2026). Role: the facts of selection, which are the
reception data. Rights: Lectionary text copyright CCD, refrains
copyright ICEL --- described by reference, never quoted.
\item \work{Liturgia Horarum}, Lauds of 15 August, Latin served at
societaslaudis.org, cross-checked at universalis.com (exact URLs in the
reception notes). Both derivative witnesses; the editio typica altera
not opened. Role: the \term{arcam incorruptibilem} intercession, Latin
with editorial paraphrase.
\item Byzantine service texts, described rather than quoted because the
accessible English translations are unidentified and presumptively
copyrighted: the Reader Menaion for 15 August at st-sergius.org
(Emenaion 08-15.pdf); the Antiochian Archdiocese Great Vespers text for
14 August; the Greek Orthodox Archdiocese feast and service
descriptions at goarch.org (the Entrance readings; the consecration
dialogue). Role: the Dormition psalmody and canon (Ps 131:8, 11; the
animate ark; Jephonias; the Ode 6 catalogue), the Entrance lections
(with their Protoevangelium substrate disclosed), the entrance liturgy.
\item Akathist hymn, ecclesiastical Greek read at Greek
Wikisource (the \term{Akathistos hymnos} page; exact URL in the leaf's
source audit) and in the transcription at lyricstranslate.com. Role: stanza 23,
\term{kib\=ote chrys\=otheisa t\=o Pneumati}; anonymous, fifth--sixth
century; the Romanos ascription unproven. Original PD; renderings
editorial.
\item CANTUS indices: Cantus Index
(\url{https://cantusindex.org/id/006069a} and correspondingly numbered
pages) and the CANTUS Database full-text search (cantusdatabase.org),
retrieved August 2026. Role: the medieval-office negatives and the late
positives (\work{Arca dei in qua reconditur}; \work{Habitavit arca
domini}, text truncated, manuscript pending).
\end{itemize}

\subsection*{Patristic and medieval witnesses}

Identified public-domain translations are named at every quotation;
texts with no public-domain English are given in the original language
with labeled editorial paraphrases, per the rights rules.

\begin{itemize}
\item \work{Ante-Nicene Fathers} (Roberts--Donaldson American edition;
public domain), New Advent, CCEL, and Early Christian Writings
transcriptions used as delivered: Justin, \work{Dialogue with Trypho}
36, 85 (vol.~1); Hippolytus, exegetical fragments (vol.~5,
\url{https://www.newadvent.org/fathers/0502.htm}), the Daniel fragment
quoted with its translation artifact corrected; Ps.-Gregory
Thaumaturgus, Annunciation homilies (vol.~6, ccel.org), with the
editors' own spuria note; Ps.-Methodius, \work{De Simeone et Anna}
(vol.~6), pseudonymous, fifth--sixth century or later; Victorinus,
\work{Commentary on the Apocalypse}, Jerome's recension (vol.~7,
\url{https://www.newadvent.org/fathers/0712.htm}); \work{Gospel of
Nicodemus} (vol.~8, earlychristianwritings.com).
\item \work{Nicene and Post-Nicene Fathers} (public domain), New Advent
transcriptions: Augustine, \work{Enarrationes in Psalmos} 131, trans.\
Tweed, NPNF I/8, \url{https://www.newadvent.org/fathers/1801132.htm}
--- the Latin negative; Ephrem, Nativity hymns, NPNF II/13.
\item Ephrem, \work{Hymns on the Nativity} (Nat.~4; 16.16--17), quoted
in the wording of J.~B.~Morris, \work{Select Works of S.~Ephrem the
Syrian} (Oxford, 1847; public domain),
\url{https://archive.org/details/selectworksofsep41ephr}. McVey's 1989
translation (in copyright) was read at archive.org for her editorial
judgment that the ark of Nat.~4.113 is Mary; Brock's guide
(\url{https://syri.ac/brock/ephrem}) is the genuineness control; the
Ephraem Graecus prayers are excluded as pseudonymous.
\item Hippolytus, \work{Commentary on Daniel} IV.24, trans.\
T.~C.~Schmidt (2010), witness PDF cited in the notes (pergrazia.com).
In copyright; brief attributed use. Role: the content-match giving the
ark to Christ's own body.
\item \work{Patrologia Graeca} (Migne; public domain), archive.org
scans (\url{https://archive.org/details/PatrologiaGraeca}): Proclus,
Or.~V, PG 65:719--720 (genuine per Constas, reported) and Or.~VII, PG
65:759--760 (authenticity unsorted); Hesychius, Homily V \work{De
S.~Maria Deipara}, PG 93:1461--1464 (corrected columns; Aubineau's
disposition pending); the \work{Collectio dictorum} among Cyril of
Alexandria's works, PG 77:1221--1224, doubtful and so labeled; Cyril,
Ephesus Homily 4, PG 77:991--996, and Homily 11, PG 77:1029--1040 ---
verified negatives for \term{kib\=otos}. Greek and Latin quoted with
labeled paraphrases; short English phrases credited to Livius.
\item John of Damascus, Dormition homilies, PG 96:700--761 (Hom.~II at
722--754, reported columns). English: the circulating text of Hom.~II
\S 2, read at malankaraworld.com, appears to be Mary H.~Allies (1898;
public domain), the identification not confirmed line by line and so
flagged; the Latin verified in the 1942 Breviarium. Role: the summit
text and Roman Matins lesson 1568--1950.
\item Chrysippus of Jerusalem, \work{Oratio in sanctam Mariam
Deiparam}, ed.\ Jugie, Patrologia Orientalis 19.3 (1926) 336--343,
scan at \url{https://archive.org/details/patrologiaorient19pariuoft};
PO 19 is US-public-domain; English snippets Livius's. Setting per the
Oxford Cult of Saints record E04681 (URL in the notes; the Hesychius
record E05888 reported).
\item The doublet sermon, under its own honest label: an anonymous
Latin sermon on David's dance and the ark, transmitted under both
Ambrose's name (PL 17:689, among the spuria) and Maximus of Turin's
(PL 57:737--740), probably fifth--sixth century or later; Latin read in
the PL 57 scan cited in the notes, with Bruni's admonitio; English
excerpts Livius p.~77; Mutzenbecher's disposition pending.
\item Jacob of Serugh, homilies on the Annunciation and Visitation,
ed.\ Bedjan (Paris, 1902), pp.~614--774; English in M.~Hansbury,
\work{On the Mother of God} (1998; in copyright), reported with loci
via J.~Puthuparampil, \work{Mariological Thought of Mar Jacob of
Serugh}, read in full at
\url{https://archive.org/details/mariologicalthou0000jame}.
Paraphrased; no long quotation pending re-verification against
Hansbury.
\item The Turin homily attributed to Athanasius: ed.\ Lefort, \work{Le
Mus\'eon} 71 (1958). No PD English; described, not quoted; attribution
disputed, no located scholarly defense.
\item Gregory the Great, \work{Homiliae in Evangelia} 25.3 (PL
76:1189--1196), Latin read at la.wikisource.org; labeled paraphrase.
\item \work{Patrologia Latina} scans for the Latin negatives:
\work{Cogitis me} (Ps.-Jerome; Paschasius Radbertus), PL 30:121--147
--- verified negative for \term{arca}, OCR-grade caveat; Bernard,
Assumption sermons, PL 183 --- the sole \term{arca} Noah's (archive.org
items in the notes). Ambrose, \work{De institutione virginis} 8.52 (PL
16), reported for the \term{porta clausa} figure, the column not
independently verified; Jerome's parallel (\work{Comm.\ in
Hiezechielem} XIII) likewise reported.
\item Irenaeus, \work{Demonstration of the Apostolic Preaching} 84,
read at tertullian.org; paraphrase only.
\item The ancient Apocalypse commentators, in the Fathers of the Church
series, read in the archive.org collection cited in the notes:
Oecumenius, trans.\ Suggit, FOTC 112 (2006); Andrew of Caesarea,
trans.\ Constantinou, FOTC 123 (2011); Tyconius, \work{Expositio
Apocalypseos} (CCSL 107A, ed.\ Gryson), trans.\ Gumerlock, FOTC 134
(2017); Cassiodorus and the anonymous Greek scholia, FOTC 144 (2022).
In copyright; quotations at brief attributed scale, the rest
paraphrase. Role: the commentary audit at Apoc 11:19--12:1.
\item Primasius, \work{Commentarius in Apocalypsin} (CCSL 92, ed.\
Adams), Latin verified against the full transcription at
turretinfan.blogspot.com (August 2023 post), carried at transcription
level with that caution; paraphrases editorial.
\item Bede, \work{The Explanation of the Apocalypse}, trans.\
E.~Marshall (Oxford, 1878; public domain),
\url{https://archive.org/details/explanationapoc00bedegoog}; quoted
verbatim.
\item Thomas Aquinas, \work{Summa theologiae} I-II q.~102 a.~4 ad 6,
English Dominican translation (public domain),
\url{https://www.newadvent.org/summa/2102.htm}; the negative across ST
III qq.~27--37 checked on the same site.
\end{itemize}

\subsection*{Historical scholarship}

\begin{itemize}
\item Frederick Field, \work{Origenis Hexaplorum quae supersunt}
(Oxford, 1875), vol.~I, page images at
\url{https://archive.org/details/origenhexapla01unknuoft}. Public
domain. Role: Symmachus's \term{skirt\=onta} (p.~555 n.~28); the Old
Greek \term{vacat} at 2 Kings 6:7; the threshing-floor variants.
\item J.~Gildemeister, ed., Theodosius, \work{De situ terrae sanctae}
(1882), \url{https://archive.org/details/theodosiusdesit00theogoog}.
Public domain. Role: the oldest Ein Karem witness (\S 35, a distance
only); renderings editorial.
\item Thomas Livius, \work{The Blessed Virgin in the Fathers of the
First Six Centuries} (London: Burns \& Oates, 1893). Public domain.
Role: the identified PD English for the Hesychius and Chrysippus
snippets and the doublet-sermon excerpts. Limit: he cites Migne
appendices without flagging spuria; two propagated citation errors are
corrected in the notes.
\item A.~De Santi, ``Litany of Loreto,'' \work{The Catholic
Encyclopedia}, vol.~9 (1913), read at Wikisource. Public domain. Role:
the litany's history --- the 1558 print, the 1587 approbation, the
1601 fixing.
\item Alphonsus Liguori, \work{The Glories of Mary} (1750), first
complete English translation (New York, 1862; 1888 printing), searched
in full at \url{https://archive.org/details/thegloriesofmary00liguuoft}.
Public domain. Role: the devotional assembly of the Visitation
``living ark'' and the Ps 131:8--2 Kings 6 fusion.
\item Cornelius a Lapide, \work{The Great Commentary}, trans.\
Mossman, and the Latin \work{Commentaria}, searched at the witnesses
in the notes (catholicapologetics.info; archive.org). Role: the
verified negative at Luke 1:39--56. Limits: Mossman abridges; the
Latin OCR is poor; an Antwerp page-image check remains open.
\item John Henry Newman, \work{Meditations and Devotions},
\work{Discourses to Mixed Congregations} 17--18, the searched sections
of the \work{Letter to Pusey}, from newmanreader.org. Public domain.
Role: the eloquent negative --- zero occurrences of ``ark.''
\item Matthias Scheeben, \work{Mariology} I, trans.\ Geukers (1946;
German 1882), archive.org scan in the notes. Translation in copyright;
brief attributed quotations. Role: the ark in modern dogmatics, with
Scheeben's own Christological policing.
\item Jean-Baptiste Chautard, \work{The Soul of the Apostolate}
(Gethsemani translation, 1946), read at the mountsaintbernard.org PDF.
Copyright unresolved; paraphrased, never quoted. Role: the
cistern-and-channel and monstrance figures, restated; register model
for the coda.
\item Joseph Ratzinger, \work{Daughter Zion} (Ignatius, 1983; German
1977), quoted from the Ignatius Press excerpt in \work{Catholic World
Report} (25 March 2025; URL in the notes); page numbers not
established. In copyright; brief attributed quotations.
\item Marie E.~Isaacs, ``Mary in the Lucan Infancy Narrative,''
\work{The Way Supplement} 25 (1975),
\url{https://www.theway.org.uk/back/s025Isaacs.pdf}, read in full.
Role: the critical school's soberest statement, quoted at her own
grounds. In copyright; brief attributed quotations.
\item Jan M.~Koz{\l}owski, ``Mary as the Ark of the Covenant in the
Scene of the Visitation (Luke 1:39--56) Reconsidered,''
\work{Warszawskie Studia Teologiczne} 31/1 (2018) 108--116,
\url{https://czasopismowst.pl/index.php/wst/article/view/159}, read in
full. Role: the modern revival; also the transmitting witness for the
verbatim quotations of Laurentin, Brown, Fitzmyer, Sch\"urmann, Bovon,
Strauss, and Rice (his notes 8--10, checked in full). His 2017
\work{ETL} article and 2022 \work{Antigone} essay are reported.
\item Jacobus d.~W. de Koning, ``Maria as anti-tipe van die
verbondsark in Lukas 1:39--56,'' \work{In die Skriflig} 54(1), a2561
(2020),
\url{https://indieskriflig.org.za/index.php/skriflig/article/view/2561},
read in full. Role: the Reformed statement of the typology.
\item The Lucan debate, consulted at one remove (loci as transmitted
in Koz{\l}owski 2018 and Isaacs, or from publishers' records; none
opened at its own pages): Burrows, \work{The Gospel of the Infancy}
(1940); Lyonnet, \work{Biblica} 20 (1939); Laurentin, \work{Structure
et th\'eologie de Luc I--II} (1957) and \work{A Short Treatise on the
Virgin Mary} (ET 1991); Brown, \work{The Birth of the Messiah} (1993);
Fitzmyer (AB 28, 1981); Sch\"urmann (1969); Bovon (1989; Hermeneia
2002; \work{Luke the Theologian}); Strauss (1995); Rice (2016);
Goulder and Sanderson, \work{JTS} 8 (1957); Serra, \work{La Donna
dell'Alleanza} (2006); Feuillet, \work{J\'esus et sa m\`ere} (1974)
and \work{RB} 66 (1959); de la Potterie, \work{Mary in the Mystery of
the Covenant} (1992).
\item Ark-and-cherubim scholarship, consulted at one remove (reviews,
surveys, publishers' listings): de Vaux, \work{Ancient Israel} (ET
1961); Haran (1978); Mettinger (1982); Seow (ABD 1, 1992); von Rad (ET
1966); Eichler, \work{The Ark and the Cherubim} (2021); Ahlstr\"om
(\work{JNES} 43, 1984); van der Toorn and Houtman (\work{JBL} 113,
1994); Kitchen (2003); Noegel (2015); Miller and Roberts, \work{The
Hand of the Lord} (1977); Rost (1926); McCarter (AB 9, as reported);
with Gunkel and Mowinckel for the entrance-liturgy hypothesis, and
Lightfoot's \work{Horae Hebraicae} as transmitted by Adam Clarke.
\item Site archaeology: the Shiloh sequence and burn layer, read at
the Emek Shaveh and Associates for Biblical Research pages cited in
the notes, the excavators' claims reported at one remove; the Jericho
literature (Sellin--Watzinger, Garstang, Kenyon) and Zertal's Ebal
claims, reported, not re-examined.
\item Patristic scholarship, reported: Constas, \work{Proclus of
Constantinople and the Cult of the Virgin in Late Antiquity};
Aubineau, \work{Les hom\'elies festales d'H\'esychius de J\'erusalem}
(not yet consulted --- open item, section 90); Mutzenbecher's CCSL 23
(pending); the Oxford Cult of Saints records named above; DiPippo's
compendium of breviary reforms at newliturgicalmovement.org.
\item Apocalypse scholarship, reported at one remove: Bauckham,
\work{The Climax of Prophecy} (1993), via a transcription at
exegesisandtheology.com; Aune (WBC 52B, 1998) and Beale (NIGTC, 1999)
as reported at readingacts.com. None consulted in the editions named.
\end{itemize}

\subsection*{Jewish, Samaritan, Islamic, and other ancient sources}

\begin{itemize}
\item Josephus, \work{Antiquities} and \work{Jewish War}, in William
Whiston's translation (public domain), Project Gutenberg ebooks 2848
and 2850 (\url{https://www.gutenberg.org/ebooks/2848};
\url{https://www.gutenberg.org/ebooks/2850}), with Ant.\ 18.85--89
verified at lexundria.com and
the Greek at the Scaife Viewer (URNs in the notes). Role: the seventy
of Bethsames with Whiston's note; the Samaritan sacred-vessels episode
of AD 36; the empty Second-Temple sanctuary.
\item \work{Vitae Prophetarum}, Jeremias chapter, ed.\ and trans.\
C.~C.~Torrey (SBL Monograph 1, 1946), read at
\url{http://lotp.thebookofenoch.info/} with the edition scan at
\url{https://archive.org/details/SBLMS1}; Hare's OTP rendering
reported. The 1946 renewal status was not established, so the chapter
is paraphrased; the possibly Christian sign-clause is flagged in use.
\item 2 Baruch 6, trans.\ R.~H.~Charles (1896; public domain),
verified at
\url{https://www.pseudepigrapha.com/pseudepigrapha/2Baruch.html},
cross-checked at the Wesley Center Online; the ``holy ark''/``ephod''
caution printed with the quotation.
\item 4 Baruch 3, in the public-domain translation of Kraft and
Purintun, verified at
\url{http://ccat.sas.upenn.edu/rak//publics/pseudepig/ParJer-Eng.html}.
\item 4 Esdras 10:21--22, described from the RSV at Bible Gateway per
the rights rules; an appendix witness, the dissenting ``plundered''
voice.
\item Mishnah and Babylonian Talmud: m.\ Sheqalim 6:1--2; b.\ Yoma 52b,
53b--54a, read via the Sefaria API
(\url{https://www.sefaria.org/Yoma.52b} and correspondingly named
pages); b.\ Rosh Hashanah 31a with Lamentations Rabbah, proem 25,
reported. Rights: the William Davidson (Steinsaltz) English is CC
BY-NC and is not quoted; all rabbinica paraphrased with exact loci.
\item Tacitus, \work{Historiae} 5.9, Latin verified at
\url{https://www.thelatinlibrary.com/tacitus/tac.hist5.shtml}; English
editorial. Role: the pagan witness to the empty sanctuary.
\item Qur'an 2:248, in Pickthall, \work{The Meaning of the Glorious
Koran} (1930; US public domain since 2026), verified via the quran.com
API (\url{https://quran.com/2/248}); the tafsir elaborations and Mahdi
traditions are reported leads, loci unverified.
\item \Needspace{14\baselineskip}\work{Kebra Nagast} (early fourteenth century; Yeshaq of Aksum),
dated on the colophon as fixed by Conti Rossini, Littmann, and Cerulli
(reported); verified at the level of the sourced survey at
\url{https://en.wikipedia.org/wiki/Kebra_Nagast} and Budge's 1922
introduction at sacred-texts.com; Ullendorff's 1941 examination
reported via Munro-Hay. Role: reception history of the Aksum claim and
the \term{tabot} practice.
\end{itemize}
